% GENERATED FILE. Edit canonical lessons, then run npm run generate:books.
% canonical-source-hash: fnv1a64:0a997312987b901f
% canonical-lessons: IT-C290-crimine, IT-C290-delitto, IT-C290-criminale, IT-C290-vittima, IT-C290-commettere

\chapter{Delitti e vittime}
\label{ch:it-delitti-e-vittime}
\hlchaptermodality{290}

\begin{chapteropening}
\textbf{By the end of this chapter:} \emph{I can talk about a crime, the criminal and the victim, and say what someone committed.}

Complete the last lesson of chapter 290: I can talk about a crime, the criminal and the victim, and say what someone committed.
\end{chapteropening}

\section[il crimine]{il crimine — the crime}

\label{lesson:IT-C290-crimine}



\begin{quote}
Before the new one: say the Italian for the audience, then the Italian for the editorial.
\end{quote}

\subsection*{You'll want to know: il crimine}
\textbf{il crimine} — \textquotedblleft{}the crime\textquotedblright{}.

\begin{quote}
È un crimine molto grave.
\end{quote}

\begin{quote}
\emph{It is a very serious crime.}
\end{quote}

\emph{La lotta al crimine} — the fight against crime.

\begin{cousinweb}
\emph{crimine} is inherited from Latin \emph{crīmen} \textquotedblleft{}a charge, a crime\textquotedblright{} (EtymDB).
\end{cousinweb}

\begin{grammarlens}[title={What you've built}]
One more word for telling, explaining and describing at B1.
\end{grammarlens}

\subsection*{Your turn}
\begin{itemize}
\raggedright
  \item \emph{Say it:} \emph{il crimine}
  \item \emph{Say it:} \emph{il crimine}, once more
  \item \emph{Say it:} a sentence of your own with \emph{il crimine}
  \item \emph{Recall:} say the Italian for the audience, then the Italian for the editorial, then say \emph{il crimine} again
\end{itemize}

\begin{tcolorbox}[breakable,colback=teal!4,colframe=teal!35!black,title={Before you move on}]
What does il crimine mean? (\textquotedblleft{}The crime\textquotedblright{}.) Say it once more.
\end{tcolorbox}
\section[il delitto]{il delitto — the serious crime; the murder}

\label{lesson:IT-C290-delitto}



\begin{quote}
Before the new one: say the Italian for the editorial, then the Italian for the crime.
\end{quote}

\subsection*{You'll want to know: il delitto}
\textbf{il delitto} — \textquotedblleft{}the serious crime; the murder\textquotedblright{}.

\begin{quote}
Il delitto è avvenuto in una villa.
\end{quote}

\begin{quote}
\emph{The crime took place in a villa.}
\end{quote}

In the papers \emph{un delitto} is often a killing.

\begin{cousinweb}
\emph{delitto} is a learned borrowing from Latin \emph{dēlictum} \textquotedblleft{}a fault, a misdeed\textquotedblright{} (EtymDB), from \emph{dēlinquere} \textquotedblleft{}to fall short, to do wrong\textquotedblright{}.
\end{cousinweb}

\begin{grammarlens}[title={What you've built}]
One more word for telling, explaining and describing at B1.
\end{grammarlens}

\subsection*{Your turn}
\begin{itemize}
\raggedright
  \item \emph{Say it:} \emph{il delitto}
  \item \emph{Say it:} \emph{il delitto}, once more
  \item \emph{Say it:} a sentence of your own with \emph{il delitto}
  \item \emph{Recall:} say the Italian for the editorial, then the Italian for the crime, then say \emph{il delitto} again
\end{itemize}

\begin{tcolorbox}[breakable,colback=teal!4,colframe=teal!35!black,title={Before you move on}]
What does il delitto mean? (\textquotedblleft{}The serious crime; the murder\textquotedblright{}.) Say it once more.
\end{tcolorbox}
\section[il criminale]{il criminale — the criminal}

\label{lesson:IT-C290-criminale}



\begin{quote}
Before the new one: say the Italian for the crime, then the Italian for the serious crime.
\end{quote}

\subsection*{You'll want to know: il criminale}
\textbf{il criminale} — \textquotedblleft{}the criminal\textquotedblright{}.

\begin{quote}
Il criminale è stato visto vicino alla stazione.
\end{quote}

\begin{quote}
\emph{The criminal was seen near the station.}
\end{quote}

Also an adjective: \emph{un'azione criminale}.

\begin{cousinweb}
\emph{criminale} is a learned borrowing from Late Latin \emph{crīminālis} \textquotedblleft{}of a crime\textquotedblright{}, built on \emph{crīmen} (Lewis \& Short).
\end{cousinweb}

\begin{grammarlens}[title={What you've built}]
One more word for telling, explaining and describing at B1.
\end{grammarlens}

\subsection*{Your turn}
\begin{itemize}
\raggedright
  \item \emph{Say it:} \emph{il criminale}
  \item \emph{Say it:} \emph{il criminale}, once more
  \item \emph{Say it:} a sentence of your own with \emph{il criminale}
  \item \emph{Recall:} say the Italian for the crime, then the Italian for the serious crime, then say \emph{il criminale} again
\end{itemize}

\begin{tcolorbox}[breakable,colback=teal!4,colframe=teal!35!black,title={Before you move on}]
What does il criminale mean? (\textquotedblleft{}The criminal\textquotedblright{}.) Say it once more.
\end{tcolorbox}
\section[la vittima]{la vittima — the victim}

\label{lesson:IT-C290-vittima}



\begin{quote}
Before the new one: say the Italian for the serious crime, then the Italian for the criminal.
\end{quote}

\subsection*{You'll want to know: la vittima}
\textbf{la vittima} — \textquotedblleft{}the victim\textquotedblright{}.

\begin{quote}
La vittima è un uomo di cinquant'anni.
\end{quote}

\begin{quote}
\emph{The victim is a fifty-year-old man.}
\end{quote}

Feminine in grammar, whoever the person is.

\begin{cousinweb}
\emph{vittima} is a learned borrowing from Latin \emph{victima} \textquotedblleft{}a beast for sacrifice\textquotedblright{} (EtymDB).
\end{cousinweb}

\begin{grammarlens}[title={What you've built}]
One more word for telling, explaining and describing at B1.
\end{grammarlens}

\subsection*{Your turn}
\begin{itemize}
\raggedright
  \item \emph{Say it:} \emph{la vittima}
  \item \emph{Say it:} \emph{la vittima}, once more
  \item \emph{Say it:} a sentence of your own with \emph{la vittima}
  \item \emph{Recall:} say the Italian for the serious crime, then the Italian for the criminal, then say \emph{la vittima} again
\end{itemize}

\begin{tcolorbox}[breakable,colback=teal!4,colframe=teal!35!black,title={Before you move on}]
What does la vittima mean? (\textquotedblleft{}The victim\textquotedblright{}.) Say it once more.
\end{tcolorbox}
\section[commettere]{commettere — to commit (a crime, an error)}

\label{lesson:IT-C290-commettere}



\begin{quote}
Before the new one: say the Italian for the criminal, then the Italian for the victim.
\end{quote}

\subsection*{You'll want to know: commettere}
\textbf{commettere} — \textquotedblleft{}to commit (a crime, an error)\textquotedblright{}.

\begin{quote}
Ha commesso un errore grave.
\end{quote}

\begin{quote}
\emph{He made a serious mistake.}
\end{quote}

Like \emph{mettere}: \emph{ha commesso}.

\begin{cousinweb}
\emph{commettere} is inherited from Latin \emph{committere} \textquotedblleft{}to bring together; to commit\textquotedblright{} (EtymDB), from \emph{con-} and \emph{mittere} \textquotedblleft{}to send\textquotedblright{}.
\end{cousinweb}

\begin{grammarlens}[title={What you've built}]
One more word for telling, explaining and describing at B1.
\end{grammarlens}

\subsection*{Your turn}
\begin{itemize}
\raggedright
  \item \emph{Say it:} \emph{commettere}
  \item \emph{Say it:} \emph{commettere}, once more
  \item \emph{Say it:} a sentence of your own with \emph{commettere}
  \item \emph{Recall:} say the Italian for the criminal, then the Italian for the victim, then say \emph{commettere} again
\end{itemize}

\begin{tcolorbox}[breakable,colback=teal!4,colframe=teal!35!black,title={Before you move on}]
What does commettere mean? (\textquotedblleft{}To commit (a crime, an error)\textquotedblright{}.) Say it once more.
\end{tcolorbox}
